\section{Dependencies and normalization}
\label{app:v94-dependencies}
The present article directly loads the three new modules 223--225.
The complete technical archive contains all 222 inherited core files
and the full preceding article. The principal dependencies are as
follows; module numbers refer to their unchanged source files.

The exact event, guard and arithmetic reference are the retained
physical problem. Modules 174, 177, 180--182 and 196 provide the
endpoint field, contact motion, optical recurrence, full angular
density and Borel segment parametrization. Module 220 gives the
absolute optical cancellation, complete-sheet traces and fixed-count
variation. Module 221 gives the exact signed prefix identity used in
\eqref{eq:v94-signed-input}. These ingredients retain the same mixed
endpoint-angle and internal-arclength convention.

The backward-end estimate is the two-sided physical-end theorem of
module 211; the precise display used here is restated as \eqref{eq:v94-end-finite}. The small-gradient volume estimate is
obtained from the actual fractional insertion in module 187, with its
finite-count remainder as restated in the inherited journal input
\texttt{eq:v83-fractional-finite-input}. The positive-cutoff layer cake
was audited in module 196 and used for complete noncritical growth in
module 222. Its lower cutoff is fixed before the collision limit in
every application here. These dependencies are explicitly restated in
Inputs~\ref{input:v94-volume} and \ref{input:v94-scalar}; no surface
spectral insertion is assumed.

The first new theorem uses the incoming, rather than outgoing, trace
and the actual backward layer. The second theorem uses the common
source on a target-dependent radial shell. The third uses the BV
Stieltjes identity for finite signed roof measures \cite{AFP}; it is
proved directly to fix the ramp direction and the sign of all cuts.
No Jordan decomposition is interchanged with projection. Historical
fixed-count definability and threshold integration remain credited to
\cite{Coste,CMInt}; neither gives a collision-normalized component count.

All interval lengths and cutoffs in the quantitative estimates are
fixed before their collision limits. The illustrative choice
$r(B)=c\varepsilon(B)^2$ also fixes the shell width $\sqrt{r(B)}$
first. It estimates the displayed radial-transit terms but does not
remove either remaining concentration. The two revision-88 source
variants remain separately preserved; their incompatible module
histories are not merged.
